% GENERATED FILE. Edit canonical lessons, then run npm run generate:books.
% canonical-source-hash: fnv1a64:266f2a1e312f345a
% canonical-lessons: ZH-C196-chepiao, ZH-W196-piao-ticket, ZH-C196-jipiao, ZH-C196-hangban, ZH-C196-qifei, ZH-C196-daoda

\chapter{Tickets and Flights}
\label{ch:zh-tickets-and-flights}
\hlchaptermodality{196}

\begin{chapteropening}
\textbf{By the end of this chapter:} \emph{I can buy a train ticket and a plane ticket, and say when a flight takes off and when a train arrives.}

Complete the last lesson of chapter 196: I can buy a train ticket and a plane ticket, and say when a flight takes off and when a train arrives.
\end{chapteropening}

\section[\texorpdfstring{\zh{车票} (chēpiào)}{车票 (chēpiào)}]{\zh{车票} (chēpiào) — train or bus ticket}

\label{lesson:ZH-C196-chepiao}



\begin{quote}
Before the new one: say the Mandarin for password, then the Mandarin for bank account.
\end{quote}

\subsection*{You'll want to know: \zh{车票}}
\textbf{\zh{车票}} — \emph{chēpiào} — \textquotedblleft{}train or bus ticket\textquotedblright{}.

\begin{quote}
\textbf{\zh{车票多少钱}} \emph{chēpiào duōshao qián} — How much is the ticket?
\end{quote}

The new character, in pieces: \zh{票} — eleven strokes: on top a bar over a frame with two short verticals; under it \zh{二}, a hooked vertical and two dots. The next lesson puts it in your hand.

\begin{cousinweb}
\textbf{\zh{车票}} (\emph{chēpiào}) is built from \zh{车} 'vehicle' and \zh{票} 'a ticket'. The first is inherited from Middle Chinese \emph{tsyhae} (Guangyun, in Baxter's transcription); no dictionary we use records an older form of the second, so its history is not settled.
\end{cousinweb}

\begin{grammarlens}[title={What you've built}]
One more word for reading notices and messages.
\end{grammarlens}

\subsection*{Your turn}
\begin{itemize}
\raggedright
  \item \emph{Say it:} \emph{chēpiào}
  \item \emph{Say it:} \emph{chēpiào}, once more
  \item \emph{Say it:} \emph{chēpiào duōshao qián} — How much is the ticket?
  \item \emph{Recall:} say the Mandarin for password, then the Mandarin for bank account, then say \emph{chēpiào} again
\end{itemize}

\begin{tcolorbox}[breakable,colback=teal!4,colframe=teal!35!black,title={Before you move on}]
What does \zh{车票} mean? (\textquotedblleft{}Train or bus ticket\textquotedblright{}.) Say it once more.
\end{tcolorbox}
\section[\texorpdfstring{\zh{票} (piào)}{票 (piào)}]{\zh{票} — a framed top over \zh{二} and three strokes}

\label{lesson:ZH-W196-piao-ticket}



\begin{quote}
The word first, then the character inside it.

\begin{itemize}
\raggedright
  \item \emph{Recall:} say \emph{chēpiào} — \textbf{R1}, one lesson back, before it decays
\end{itemize}

\begin{quote}
\zh{车} \zh{二}
\end{quote}

\emph{Read it:} each one aloud, with nothing said first — \emph{chē}, \emph{èr}
\end{quote}

\subsection*{Script — a framed top over \zh{二} and three strokes}
\hlblockfigure{\detokenize{figures/ZH-W196-piao-ticket-filmstrip.pdf}}{How \zh{票} is written, stroke by stroke}

\begin{quote}
\zh{票}
\end{quote}

\textbf{Two pieces, top before bottom.}

\textbf{The top piece, six strokes.} A bar across the top. Then a frame under it: a short vertical on the left, a stroke along the top that turns down the right side, two short verticals down through the top of the frame, and a bar that closes it.

\textbf{The bottom piece, five strokes.} \zh{二}: a short bar, then a longer one. A vertical down the middle that hooks to the left at its foot. A dot falling to the left, then a dot falling to the right.

6 + 5 = eleven.

\subsection*{Your turn}
\begin{itemize}
\raggedright
  \item \emph{Write it:} \zh{票}, all eleven, with the model in front of you
  \item \emph{Cover:} write it again from memory, and count the pieces, not the strokes
\end{itemize}

\begin{tcolorbox}[breakable,colback=teal!4,colframe=teal!35!black,title={Before you move on}]
How many strokes in \zh{票}? (\textbf{Eleven}.) What starts the bottom piece? (\textbf{\zh{二}}.) What sits either side of the hooked vertical? (\textbf{Two dots}.)

Source: \href{https://raw.githubusercontent.com/chanind/hanzi-writer-data/68d10a4b21150cae5e1ebbd223eed289cf32d90c/data/\%E7\%A5\%A8.json}{Hanzi Writer Data \zh{票}, pinned stroke paths and medians}.
\end{tcolorbox}
\section[\texorpdfstring{\zh{机票} (jīpiào)}{机票 (jīpiào)}]{\zh{机票} (jīpiào) — plane ticket}

\label{lesson:ZH-C196-jipiao}



\begin{quote}
Before the new one: say the Mandarin for bank account, then the Mandarin for train or bus ticket.
\end{quote}

\subsection*{You'll want to know: \zh{机票}}
\textbf{\zh{机票}} — \emph{jīpiào} — \textquotedblleft{}plane ticket\textquotedblright{}.

\begin{quote}
\textbf{\zh{我买了十号的机票}} \emph{wǒ mǎile shí hào de jīpiào} — I've bought a plane ticket for the tenth.
\end{quote}

\begin{cousinweb}
\textbf{\zh{机票}} (\emph{jīpiào}) is built from \zh{机}, short for \emph{fēijī} 'aeroplane', and \zh{票} 'a ticket'. The first is inherited from Middle Chinese \emph{kj+j} (Guangyun, in Baxter's transcription); no dictionary we use records an older form of \zh{票}, so its history is not settled.
\end{cousinweb}

\begin{grammarlens}[title={What you've built}]
One more word for reading notices and messages.
\end{grammarlens}

\subsection*{Your turn}
\begin{itemize}
\raggedright
  \item \emph{Say it:} \emph{jīpiào}
  \item \emph{Say it:} \emph{jīpiào}, once more
  \item \emph{Say it:} \emph{wǒ mǎile shí hào de jīpiào} — I've bought a plane ticket for the tenth.
  \item \emph{Recall:} say the Mandarin for bank account, then the Mandarin for train or bus ticket, then say \emph{jīpiào} again
\end{itemize}

\begin{tcolorbox}[breakable,colback=teal!4,colframe=teal!35!black,title={Before you move on}]
What does \zh{机票} mean? (\textquotedblleft{}Plane ticket\textquotedblright{}.) Say it once more.
\end{tcolorbox}
\section[hángbān]{hángbān — scheduled flight}

\label{lesson:ZH-C196-hangban}



\begin{quote}
Before the new one: say the Mandarin for train or bus ticket, then the Mandarin for plane ticket.
\end{quote}

\subsection*{You'll want to know: hángbān}
\emph{hángbān} — \textquotedblleft{}scheduled flight\textquotedblright{}.

\begin{quote}
\emph{zhège hángbān wǎndiǎn le} — This flight is delayed.
\end{quote}

Its characters come later; for now it is a word for the ear.

\begin{cousinweb}
\emph{hángbān} is built from \emph{háng} 'to sail' and \emph{bān} 'a shift, a set run', both inherited from Middle Chinese, \emph{hang} and \emph{paen} (Guangyun, in Baxter's transcription). The dictionary glosses the first 'a boat' and the second as dividing jade tokens; how the pair came to mean a flight is not settled.
\end{cousinweb}

\begin{grammarlens}[title={What you've built}]
One more word for reading notices and messages.
\end{grammarlens}

\subsection*{Your turn}
\begin{itemize}
\raggedright
  \item \emph{Say it:} \emph{hángbān}
  \item \emph{Say it:} \emph{hángbān}, once more
  \item \emph{Say it:} \emph{zhège hángbān wǎndiǎn le} — This flight is delayed.
  \item \emph{Recall:} say the Mandarin for train or bus ticket, then the Mandarin for plane ticket, then say \emph{hángbān} again
\end{itemize}

\begin{tcolorbox}[breakable,colback=teal!4,colframe=teal!35!black,title={Before you move on}]
What does hángbān mean? (\textquotedblleft{}Scheduled flight\textquotedblright{}.) Say it once more.
\end{tcolorbox}
\section[qǐfēi]{qǐfēi — to take off (of a plane)}

\label{lesson:ZH-C196-qifei}



\begin{quote}
Before the new one: say the Mandarin for plane ticket, then the Mandarin for scheduled flight.
\end{quote}

\subsection*{You'll want to know: qǐfēi}
\emph{qǐfēi} — \textquotedblleft{}to take off (of a plane)\textquotedblright{}.

\begin{quote}
\emph{fēijī shí diǎn qǐfēi} — The plane takes off at ten.
\end{quote}

Its characters come later; for now it is a word for the ear.

\begin{cousinweb}
\emph{qǐfēi} is built from \zh{起} 'to rise' and \emph{fēi} 'to fly', both inherited from Middle Chinese, \emph{khiX} and \emph{pj+j} (Guangyun, in Baxter's transcription).
\end{cousinweb}

\begin{grammarlens}[title={What you've built}]
One more word for reading notices and messages.
\end{grammarlens}

\subsection*{Your turn}
\begin{itemize}
\raggedright
  \item \emph{Say it:} \emph{qǐfēi}
  \item \emph{Say it:} \emph{qǐfēi}, once more
  \item \emph{Say it:} \emph{fēijī shí diǎn qǐfēi} — The plane takes off at ten.
  \item \emph{Recall:} say the Mandarin for plane ticket, then the Mandarin for scheduled flight, then say \emph{qǐfēi} again
\end{itemize}

\begin{tcolorbox}[breakable,colback=teal!4,colframe=teal!35!black,title={Before you move on}]
What does qǐfēi mean? (\textquotedblleft{}To take off (of a plane)\textquotedblright{}.) Say it once more.
\end{tcolorbox}
\section[dàodá]{dàodá — to reach a destination; to arrive}

\label{lesson:ZH-C196-daoda}



\begin{quote}
Before the new one: say the Mandarin for scheduled flight, then the Mandarin for to take off.
\end{quote}

\subsection*{You'll want to know: dàodá}
\emph{dàodá} — \textquotedblleft{}to reach a destination; to arrive\textquotedblright{}.

\begin{quote}
\emph{huǒchē xiàwǔ sān diǎn dàodá} — The train arrives at three in the afternoon.
\end{quote}

Its characters come later; for now it is a word for the ear.

\begin{cousinweb}
\emph{dàodá} is built from \zh{到} 'to arrive' and \emph{dá} 'to get through', both inherited from Middle Chinese, \emph{tawH} and \emph{dat} (Guangyun, in Baxter's transcription).
\end{cousinweb}

\begin{grammarlens}[title={What you've built}]
One more word for reading notices and messages.
\end{grammarlens}

\subsection*{Your turn}
\begin{itemize}
\raggedright
  \item \emph{Say it:} \emph{dàodá}
  \item \emph{Say it:} \emph{dàodá}, once more
  \item \emph{Say it:} \emph{huǒchē xiàwǔ sān diǎn dàodá} — The train arrives at three in the afternoon.
  \item \emph{Recall:} say the Mandarin for scheduled flight, then the Mandarin for to take off, then say \emph{dàodá} again
\end{itemize}

\begin{tcolorbox}[breakable,colback=teal!4,colframe=teal!35!black,title={Before you move on}]
What does dàodá mean? (\textquotedblleft{}To reach a destination; to arrive\textquotedblright{}.) Say it once more.
\end{tcolorbox}
